\chapter{RESULT AND ANALYSIS}
\section{Environmental Sensor-Based Fire Risk Assessment}

The developed sensor prototype was tested under different environmental conditions
to verify the fire risk estimation mechanism.

The sensor node successfully collected temperature, humidity, and gas concentration
data and generated a corresponding fire risk score.

Example observations obtained from testing are shown in Table~\ref{tab:risk}.


\begin{table}[h]
\centering
\caption{Environmental Sensor Testing Results}
\label{tab:risk}
\begin{tabular}{|c|c|c|c|c|}
\hline
Condition & Temperature & Humidity & Gas Level & Risk Level\\
\hline
Normal Environment & Low & High & Low & Low\\
\hline
Warm and Dry & Medium & Medium & Medium & Moderate\\
\hline
Smoke Exposure & High & Low & High & High\\
\hline
\end{tabular}
\end{table}


The results indicate that combining multiple sensor parameters provides a more
reliable indication of fire risk compared to individual sensor threshold methods.
The sensor fusion approach reduces false alarms caused by temporary variations in
a single environmental parameter.

\section{LoRa Communication Simulation Results}

The LoRa communication simulation successfully demonstrated packet transmission
between LoRa sensor nodes using the FLORA framework.

The implemented simulation verified:

\begin{itemize}
    \item Successful packet transmission between LoRa nodes.
    \item Reception of packets at the LoRa gateway.
    \item Basic wireless communication functionality.
\end{itemize}

The simulation confirms that the FLORA and INET-based environment can be used for
evaluating the proposed forest monitoring network.

The future simulation work will include:

\begin{itemize}
    \item Implementation of LDSE multi-hop routing protocol.
    \item Energy consumption analysis.
    \item Packet delivery ratio evaluation.
    \item Network scalability testing.
    \item Performance comparison with existing routing approaches.
\end{itemize}


\section{Current Limitations}

Although the initial prototype and simulation environment have been successfully
developed, some components of the complete system are still under development.

The remaining tasks include:

\begin{itemize}
    \item Deployment of TinyML-based acoustic classification model.
    \item Implementation of LDSE routing protocol.
    \item Hardware-based LoRa communication testing.
    \item Cloud dashboard integration.
    \item Large-scale network simulation.
\end{itemize}


\section{Current Project Progress}

\begin{table}[H]
\centering
\caption{Current Project Status}
\begin{tabular}{|c|c|}
\hline
Module & Status\\
\hline
ESP32-S3 Sensor Node & Completed\\
\hline
Temperature/Humidity/Gas Sensing & Completed\\
\hline
Fire Risk Score Calculation & Completed\\
\hline
Acoustic Data Acquisition & Completed\\
\hline
TinyML Sound Classification & In Progress\\
\hline
LoRa Hardware Communication & In Progress\\
\hline
FLORA LoRa Simulation & Completed\\
\hline
LDSE Routing Protocol & Future Work\\
\hline
Web Dashboard & Future Work\\
\hline
\end{tabular}
\end{table}



\section{Acoustic Threat Detection Results}

The deployed \texttt{tiny\_int8.tflite} model was evaluated on the held-out
FSC22 test split (240 samples, stratified across the 6 project classes). The
model achieved an overall accuracy of \textbf{87.92\%} (211/240 correct) and
a micro-averaged AUC-ROC of \textbf{0.9811}, exceeding the full model's
INT8 accuracy of 86.67\% despite using roughly 7$\times$ fewer parameters
(Section~\ref{subsec:model-evolution}).

Table~\ref{tab:classification-report} reports precision, recall, and
F1-score per class. Background -- the majority class with 165 test samples
-- is classified with very high precision and recall (0.9398 / 0.9455),
which is the desired behaviour for a safety-critical detector, since false
positives on background would generate nuisance alerts.

\begin{table}[H]
    \centering
    \caption{Classification report -- Tiny INT8 model on FSC22 test split}
    \label{tab:classification-report}
    \begin{tabular}{lccccc}
        \toprule
        \textbf{Class} & \textbf{Precision} & \textbf{Recall} & \textbf{F1-score} & \textbf{Support} \\
        \midrule
        axe          & 0.7692 & 0.6667 & 0.7143 & 15 \\
        background   & 0.9398 & 0.9455 & 0.9426 & 165 \\
        chainsaw     & 0.8462 & 0.7333 & 0.7857 & 15 \\
        gunshot      & 0.6500 & 0.8667 & 0.7429 & 15 \\
        handsaw      & 1.0000 & 0.8000 & 0.8889 & 15 \\
        tree\_falling & 0.5625 & 0.6000 & 0.5806 & 15 \\
        \midrule
        \textbf{Accuracy} & \multicolumn{3}{c}{0.8792} & 240 \\
        Macro avg    & 0.7946 & 0.7687 & 0.7758 & 240 \\
        Weighted avg & 0.8853 & 0.8792 & 0.8801 & 240 \\
        \bottomrule
    \end{tabular}
\end{table}

Figure~\ref{fig:prf-chart} visualises the same precision, recall, and
F1-score values, making the weaker classes easier to spot at a glance:
handsaw's recall gap (1.000 precision vs.\ 0.800 recall), gunshot's
precision gap (0.650 precision vs.\ 0.867 recall) -- suggesting the model
is currently biased toward flagging ambiguous impulsive sounds as gunshots
rather than missing them, an acceptable trade-off for a security
application -- and tree\_falling's overall weakness across all three
metrics.

\begin{figure}[H]
    \centering
    \includegraphics[width=0.85\textwidth]{tinyml_precision_recall_f1.png}
    \caption{Per-class precision, recall, and F1-score for the deployed
    tiny model}
    \label{fig:prf-chart}
\end{figure}

Figure~\ref{fig:confusion-matrix} shows the confusion matrix underlying
these scores. Tree falling's weak F1-score (0.5806) is driven mainly by
confusion with background: five of fifteen tree-falling clips were
misclassified as background, plausibly because the low-frequency
crack/creak signature of a falling tree overlaps with wind and ambient
forest noise, and because tree falling has the smallest test support (15
samples) of any class.

\begin{figure}[H]
    \centering
    \includegraphics[width=0.72\textwidth]{tinyml_confusion_matrix.png}
    \caption{Confusion matrix for the deployed tiny INT8 model)}
    \label{fig:confusion-matrix}
\end{figure}

Despite this per-class variation, the model separates every class from the
rest well in ranking terms, achieving a micro-averaged AUC-ROC of 0.9811 and
a per-class one-vs-rest AUC no lower than 0.9384 (gunshot) -- indicating
that the confidence threshold, rather than the model's underlying
discriminative power, is the primary lever for improving the weaker classes
before final deployment.